De l'analyse de ces différents profils d'utilisateurs des collections de l'INA, une conclusion revient sans manquer, la bonne accessibilité aux collections repose en grande partie sur l'expertise des documentalistes, qui est pourtant en péril aujourd'hui. Le travail des documentalistes reste extrêmement important, que cela soit lors du traitement documentaire ou lors de la médiation de connaissances. Actuellement, si il n'y a aucune médiation des documentalistes, à l'INAthèque ou par la création de corpus, la compréhension de la base de données est compromise.

Il est donc au centre de nos préoccupations, afin de faciliter l'accessibilité aux collections de donner aux utilisateurs de nouvelles clefs de compréhension des bases de données de l'INA. Notre projet de cartographie doit donc se concentrer sur le partage des connaissances et la cartographie visuelle du traitement documentaire afin de donner toutes les clefs aux utilisateurs pour comprendre ces collections. Mais pour le mettre en place, nous devons déjà identifier quels paramètres génèrent ces difficultés à accéder à l'ensemble des collections de l'INA. Tout comme il est important de prend en compte les solutions déjà mises en place pour faciliter l'accessibilité aux collections. 